\chapter{Diseño}

\section{Visión general}

El sistema se organiza alrededor de un agente coordinador, el
Orchestrator, que media todas las interacciones entre el usuario, los
agentes especializados y la capa de persistencia compartida con la
primera fase. El cliente web mantiene un canal bidireccional con el
backend mediante WebSocket, de modo que cada paso del pipeline se refleja
en la interfaz a medida que se produce, sin necesidad de \textit{polling}.

Los diagramas de despliegue y componentes de este capítulo muestran las piezas
coordinadas por el Orchestrator: agentes especializados, entorno de simulación,
modelos cargados desde la primera fase y capa compartida de persistencia y
recuperación. El capítulo desarrolla esas piezas desde el punto de vista del
diseño de la segunda fase, no como inventario completo de todos los servicios
del proyecto.

Sobre esa estructura, el diseño de esta fase se apoya en tres decisiones
principales:

\begin{enumerate}
\item La \textbf{orquestación de agentes} del pipeline: un coordinador que
modela la sesión como una secuencia canónica de \textit{tool calls} y
reserva al modelo de lenguaje únicamente las desviaciones puntuales, en
lugar de delegarle la planificación completa.
\item La \textbf{carga dinámica de modelos}: la segunda fase descubre los
modelos aceptados por la primera, los instancia bajo demanda y los ejecuta
sin incorporarlos como dependencias estáticas del laboratorio.
\item El \textbf{laboratorio como productor y consumidor de memoria
compartida}: las observaciones de simulación se persisten en el
Knowledge Backbone y se reutilizan después para contrastar el
comportamiento observado con los postulados previos.
\end{enumerate}

Para completar esta visión con una representación formal, la
Figura~\ref{fig:despliegue-uml} muestra la arquitectura del sistema como
diagrama de despliegue UML. En ella se separan los entornos de ejecución
principales: navegador, backend de la segunda fase, servicios externos de
modelo y capa de persistencia compartida. La Figura~\ref{fig:componentes-uml}
refina esa vista con los componentes software de grano grueso que participan
en una sesión del laboratorio: la primera figura muestra dónde se ejecuta cada
pieza; la segunda baja al nivel de componentes software.

\begin{figure}[H]
\centering
\resizebox{0.92\linewidth}{!}{%
\begin{tikzpicture}[
  nodebox/.style={
    rectangle, draw=black!75, rounded corners=2pt, line width=0.6pt,
    fill=black!3, minimum height=12mm, text width=34mm,
    align=center, font=\small\sffamily
  },
  artifact/.style={
    rectangle, draw=black!60, rounded corners=1pt, line width=0.5pt,
    fill=white, minimum height=7mm, text width=29mm,
    align=center, font=\scriptsize\sffamily
  },
  storage/.style={
    cylinder, shape border rotate=90, aspect=0.25, draw=black!65,
    line width=0.55pt, fill=black!5, minimum height=10mm,
    text width=24mm, align=center, font=\scriptsize\sffamily
  }
]
  \node[nodebox] (browser) at (0,0) {<<device>>\\Navegador};
  \node[artifact] (frontend) at (0,-1.45) {Frontend\\React + Vite};

  \node[nodebox, text width=40mm] (backend) at (5.2,0) {<<execution environment>>\\Backend Phase 2};
  \node[artifact] (api) at (3.6,-1.45) {FastAPI\\REST + WS};
  \node[artifact] (orch) at (5.2,-2.65) {Orchestrator\\tools + sesión};
  \node[artifact] (sim) at (6.8,-1.45) {Runtime\\simulación 2D};

  \node[nodebox] (llm) at (10.4,0) {<<external service>>\\LLM providers};
  \node[artifact] (openrouter) at (10.4,-1.45) {OpenRouter API};

  \node[nodebox, text width=36mm] (phase1) at (0,-4.1) {<<execution environment>>\\Phase 1};
  \node[artifact] (models) at (0,-5.45) {Modelos aceptados\\código Python};

  \node[nodebox, text width=44mm] (shared) at (5.2,-4.1) {<<execution environment>>\\Infraestructura compartida};
  \node[storage] (pg) at (2.4,-6.65) {PostgreSQL\\metadatos};
  \node[storage] (minio) at (4.6,-6.65) {MinIO\\artefactos};
  \node[storage] (qdrant) at (6.8,-6.65) {Qdrant\\memorias};
  \node[storage] (neo4j) at (9.0,-6.65) {Neo4j\\grafo};

  \draw[diagarrowbi] (frontend) -- node[above, font=\scriptsize] {HTTP / WS} (api);
  \draw[diagarrow] (api) -- (orch);
  \draw[diagarrowbi] (orch) -- (sim);
  \draw[diagarrowbi] (orch) -- (openrouter);
  \draw[diagarrow] (models) |- (sim);
  \draw[diagarrowbi] (orch) -- (shared);
  \draw[diagarrow] (phase1) -- (models);
  \draw[diagarrow] (shared) -- (pg);
  \draw[diagarrow] (shared) -- (minio);
  \draw[diagarrow] (shared) -- (qdrant);
  \draw[diagarrow] (shared) -- (neo4j);
  \draw[diagarrow] (models) -- (minio);
\end{tikzpicture}}
\caption[Diagrama de despliegue UML]{Diagrama de despliegue UML del
laboratorio virtual. Se muestran los nodos de ejecución, los artefactos
desplegados y los servicios persistentes o externos utilizados durante una
sesión.}
\label{fig:despliegue-uml}
\end{figure}

\begin{figure}[H]
\centering
\resizebox{\linewidth}{!}{%
\begin{tikzpicture}[
  node distance=7mm and 8mm,
  comp/.style={
    rectangle, draw=black!75, rounded corners=2pt, line width=0.55pt,
    fill=#1, minimum height=8mm, text width=29mm,
    align=center, font=\scriptsize\sffamily
  }
]
  \node[comp=blue!7] (ui) {Cliente web\\dashboard, chat, replay};
  \node[comp=blue!4, below=of ui] (api) {API FastAPI\\REST + WebSocket};
  \node[comp=violet!9, below=of api] (orch) {Orchestrator\\estado de sesión};

  \node[comp=green!8, below left=11mm and 25mm of orch] (architect) {Architect\\especificación};
  \node[comp=green!8, below=11mm of orch] (tracker) {Tracker\\observación};
  \node[comp=green!8, below right=11mm and 25mm of orch] (analyst) {Analyst\\análisis};
  \node[comp=green!8, right=8mm of analyst] (reporter) {Reporter\\PDF};

  \node[comp=orange!9, below=15mm of tracker, text width=35mm] (runtime) {Runtime de simulación\\Environment + Agent};
  \node[comp=orange!5, left=10mm of runtime] (loader) {Model loader\\modelos fase 1};
  \node[comp=red!6, right=10mm of runtime] (charts) {Charts\\figuras del análisis};

  \node[comp=gray!8, below left=14mm and 24mm of runtime] (db) {DatabaseService\\PostgreSQL};
  \node[comp=gray!8, below=14mm of runtime] (storage) {StorageService\\MinIO};
  \node[comp=gray!8, below right=14mm and 24mm of runtime] (vector) {VectorStore\\Qdrant};
  \node[comp=gray!8, right=8mm of vector] (kg) {KnowledgeGraph\\Neo4j};
  \node[comp=gray!5, left=8mm of db] (memory) {TrackerMemoryWriter\\memoria de simulación};

  \draw[diagarrowbi] (ui) -- (api);
  \draw[diagarrowbi] (api) -- (orch);
  \draw[diagarrow] (orch) -- (architect);
  \draw[diagarrow] (orch) -- (tracker);
  \draw[diagarrow] (orch) -- (analyst);
  \draw[diagarrow] (orch) -- (reporter);
  \draw[diagarrow] (architect) -- (runtime);
  \draw[diagarrowbi] (runtime) -- (tracker);
  \draw[diagarrow] (tracker) -- (analyst);
  \draw[diagarrow] (analyst) -- (reporter);
  \draw[diagarrow] (loader) -- (runtime);
  \draw[diagarrow] (analyst) -- (charts);
  \draw[diagarrow] (charts) -- (reporter);
  \draw[diagarrow] (tracker) -- (memory);
  \draw[diagarrow] (memory) -- (db);
  \draw[diagarrow] (memory) -- (vector);
  \draw[diagarrowbi] (orch) |- (db);
  \draw[diagarrowbi] (orch) |- (storage);
  \draw[diagarrowbi] (analyst) |- (vector);
  \draw[diagarrowbi] (analyst) |- (kg);
  \draw[diagarrow] (reporter) |- (storage);
\end{tikzpicture}}
\caption[Diagrama de componentes UML]{Diagrama de componentes UML de la
segunda fase. El Orchestrator aparece como componente central y los agentes se
separan de las capas de simulación, presentación y persistencia.}
\label{fig:componentes-uml}
\end{figure}

Las secciones siguientes separan las aportaciones propias de esta fase y las
capas que materializan la sesión.

\section{Aportaciones propias de la fase}

La segunda fase parte de dos piezas que ya existen en el proyecto: los modelos
generados por la primera fase y el Knowledge Backbone compartido. El trabajo de
esta fase consiste en convertir esas piezas en un laboratorio utilizable: cargar
los modelos, ejecutarlos en un entorno común, observar lo que hacen y dejar la
sesión registrada para poder analizarla después.

Sobre esa base, la aportación principal está en la sesión de simulación. El
\textbf{Orchestrator} mantiene la conversación con el usuario y coordina a los
agentes especializados: Architect para construir el entorno, Tracker para
registrar la ejecución, Analyst para interpretar los resultados y Reporter para
generar el informe. El canal \textbf{WebSocket} hace visible ese proceso en la
interfaz, de modo que la simulación no queda reducida a una respuesta final del
modelo, sino que puede seguirse paso a paso.

La fase también añade la comparación de varios modelos en el mismo entorno,
con trazas de decisión, métricas y \textit{replay}. Esa parte es importante
porque permite observar diferencias entre modelos sin cambiar las condiciones
de la simulación. Por último, cada sesión deja persistidos sus datos principales:
conversaciones, experimentos, ejecuciones, informes y observaciones. Así, el
resultado no depende solo del estado del navegador ni de la ventana de contexto
del modelo, y las observaciones pueden volver al Knowledge Backbone como memoria
recuperable.

\section{El dominio de simulación}

La simulación se desarrolla sobre una rejilla bidimensional habitada por
agentes y recursos. Esta capa se mantiene sencilla a propósito: su único
objetivo es ofrecer un escenario común sobre el que distintos modelos de
decisión, incluso de paradigmas teóricos muy diferentes, puedan ser
evaluados en condiciones equivalentes.

El módulo \texttt{simlab.environment} define cinco conceptos: el
\texttt{Environment}, que mantiene la rejilla, el paso actual y las reglas
del entorno; el \texttt{Agent}, entidad observable situada en la rejilla; el
\texttt{Resource}, objeto con propiedades y reglas de aparición; el
\texttt{Event}, registro inmutable de cada decisión; y el
\texttt{DecisionModel}, la política de decisión cargada desde la primera fase.
La separación entre \texttt{Agent} y \texttt{DecisionModel} es deliberada:
el primero responde a \emph{quién} y \emph{dónde}, mientras que el segundo
responde a \emph{cómo decide}. Así se puede cambiar la política sin tocar la
infraestructura del entorno y ejecutar varios paradigmas en paralelo dentro de
una misma simulación sin que el entorno tenga que conocer su naturaleza.

El contrato \texttt{decide}, \texttt{update}, \texttt{get\_state}
resulta familiar para quien haya trabajado con aprendizaje por refuerzo:
\texttt{decide} resuelve la selección de acción dada una percepción, es
decir, un diccionario con la posición del agente (\texttt{x}, \texttt{y}), las
dimensiones de la rejilla (\texttt{grid\_width}, \texttt{grid\_height}),
el paso actual (\texttt{step}), los recursos visibles (\texttt{resources})
y el resultado de la última acción (\texttt{last\_action\_result});
\texttt{update} aplica el ajuste interno tras recibir la recompensa, y
\texttt{get\_state} expone el estado interno del modelo, incluida una
clave \texttt{q\_values} que el laboratorio consume para las trazas de
decisión y las gráficas de alternativas valoradas. La generalidad de esta
interfaz permite acomodar tanto modelos clásicos basados en Q-learning
como paradigmas no tabulares o estrategias deliberativas más elaboradas.

\begin{figure}[H]
\centering
\resizebox{\linewidth}{!}{%
\begin{tikzpicture}[
  node distance=8mm and 9mm,
  cls/.style={
    rectangle split, rectangle split parts=3, draw=black!75,
    rounded corners=2pt, line width=0.55pt, fill=black!3,
    text width=34mm, align=left, font=\scriptsize\sffamily
  }
]
  \node[cls] (env) {
    \textbf{Environment}
    \nodepart{second}
    width, height, step\\
    agents, resources, events
    \nodepart{third}
    add\_agent()\\
    step()\\
    run()
  };
  \node[cls, left=13mm of env] (agent) {
    \textbf{Agent}
    \nodepart{second}
    id, position\\
    model
    \nodepart{third}
    --
  };
  \node[cls, right=13mm of env] (resource) {
    \textbf{Resource}
    \nodepart{second}
    id, type\\
    position, properties
    \nodepart{third}
    --
  };
  \node[cls, below=14mm of agent] (model) {
    \textbf{DecisionModel}
    \nodepart{second}
    <<Protocol>>
    \nodepart{third}
    decide(perception)\\
    update(action, reward, perception)\\
    get\_state()
  };
  \node[cls, below=14mm of env] (event) {
    \textbf{Event}
    \nodepart{second}
    step, agent\_id\\
    action, reward\\
    perception, state
    \nodepart{third}
    --
  };
  \node[cls, below=14mm of resource, text width=38mm] (rules) {
    \textbf{ActionRule / ResourceRule}
    \nodepart{second}
    name, effect\\
    spawn policy
    \nodepart{third}
    MoveEffect\\
    ConsumeEffect\\
    NoopEffect
  };

  \node[cls, below=18mm of event, text width=37mm] (orch) {
    \textbf{Orchestrator}
    \nodepart{second}
    history, state\\
    services, tools
    \nodepart{third}
    chat()\\
    run\_simulation()\\
    analyze\_results()
  };
  \node[cls, left=12mm of orch, text width=38mm] (agents) {
    \textbf{Architect / Tracker / Analyst / Reporter}
    \nodepart{second}
    client, model
    \nodepart{third}
    run()
  };
  \node[cls, right=10mm of orch] (writer) {
    \textbf{TrackerMemoryWriter}
    \nodepart{second}
    db, vector stores
    \nodepart{third}
    write()
  };

  \draw[diagarrow] (env) -- (agent);
  \draw[diagarrow] (env) -- (resource);
  \draw[diagarrow] (env) -- (event);
  \draw[diagarrow] (agent) -- (model);
  \draw[diagarrow] (rules) -- (env);
  \draw[diagarrow] (orch) -- (agents);
  \draw[diagarrow] (orch.east) -- (writer.west);
  \draw[diagarrow] (orch.north) -- ++(0,8mm) -| (env.south);
  \draw[diagarrow] (writer.north) -- ++(0,8mm) -| (event.south east);
\end{tikzpicture}}
\caption[Diagrama de clases UML]{Diagrama de clases UML simplificado de la
segunda fase. Se incluyen las clases del dominio de simulación, el contrato
\texttt{DecisionModel} y las clases que coordinan, observan y persisten la
sesión.}
\label{fig:clases-uml}
\end{figure}

\section{El Orchestrator y el ciclo de la sesión}

El Orchestrator es la única pieza que entiende el pipeline en su
conjunto. Mantiene el historial del chat, registra el catálogo de
herramientas disponibles y despacha cada llamada al agente especializado
correspondiente. La interfaz refleja ese pipeline en el panel lateral, donde
cada agente expone su estado de ejecución y muestra el avance en tiempo real
mientras los agentes trabajan (Figura~\ref{fig:pipeline-accion}). La
planificación global no se delega en el modelo de
lenguaje, sino que se modela como una secuencia canónica de \textit{tool
calls}. Sobre esa secuencia, el LLM solo elige desviaciones puntuales:
pedir una consulta adicional al Knowledge
Backbone, repetir un paso o descartar un intento fallido.

El catálogo de herramientas combina llamadas de ejecución, como
\texttt{create\_\allowbreak environment}, \texttt{run\_\allowbreak simulation},
\texttt{analyze\_\allowbreak results} o \texttt{generate\_\allowbreak report},
con llamadas de inspección. Estas últimas no lanzan una fase nueva del
pipeline; recuperan detalle bajo demanda: episodios, trayectorias, eventos
críticos, comparaciones del Analyst o una ventana de pasos alrededor de un
instante de la simulación. Así, el Orchestrator puede responder a preguntas precisas del
usuario sin introducir de nuevo todo el JSON del Tracker en el contexto del
modelo. El catálogo completo de \textit{tool calls} expuestas a cada agente se
recoge en el Apéndice~\ref{ap:toolcalls}.

Este patrón corresponde a una composición
centralizada con organización por capas en la taxonomía reciente de
sistemas multi-agente basados en LLMs \cite{guo2024llmbased}, y se
alinea con frameworks que codifican procedimientos operativos como
secuencias de roles especializados \cite{hong2024metagpt}.

Tres responsabilidades adicionales viven exclusivamente en el
Orchestrator y no se delegan a ningún agente especializado:

\begin{enumerate}
\item \textbf{Pre-recuperación de contexto.} Antes de invocar al
Architect y al Reporter, el Orchestrator consulta el Knowledge Backbone
mediante un paso interno de pre-recuperación y entrega el resultado como
contexto de entrada; estos dos agentes no abren un canal propio al backbone.
El Analyst es la excepción: tiene \texttt{retrieve\_context} declarada como
\textit{tool call} y la invoca por su cuenta cuando necesita ampliar el
contexto durante su análisis.
\item \textbf{Persistencia diferida.} Tras recibir el JSON del Tracker,
el Orchestrator lo pasa al \texttt{TrackerMemoryWriter} para que lo
descomponga en hechos, calcule los \textit{embeddings} con Voyage en un
único \textit{batch} y los inserte de forma coordinada en PostgreSQL,
Qdrant denso y Qdrant \textit{sparse}. Si las credenciales de Voyage AI o
ZeroEntropy no están configuradas, este escritor no participa y la sesión
continúa sin escritura semántica indexada.
\item \textbf{Transmisión en vivo.} Cada salida intermedia se publica en
el canal WebSocket para que el cliente actualice la interfaz sin
esperar al final del pipeline.
\end{enumerate}

Las desviaciones que el modelo de lenguaje sí controla son de dos clases.
La primera es la \textbf{confirmación humana}: antes de lanzar una
simulación, el Orchestrator presenta al usuario la especificación de
entorno propuesta por el Architect y los modelos seleccionados, y espera
su validación; el usuario puede pedir ajustes y rehacer ese paso sin
reiniciar la sesión. La segunda es el \textbf{rerun de un paso fallido}:
si una \textit{tool call} produce una salida que no valida (por ejemplo,
un JSON de entorno malformado) el Orchestrator puede repetirla o
descartar el intento sin perder el contexto acumulado. En ambos casos la
secuencia canónica se mantiene como base; lo que el LLM decide es
\emph{cuándo} retroceder o reintentar, no el orden global del pipeline.
Esta combinación de plan fijo y desviaciones acotadas es lo que hace que
una sesión sea a la vez previsible en coste y tolerante a los fallos
puntuales propios de las salidas generativas.

\begin{figure}[p]
\centering
\resizebox{0.82\linewidth}{!}{%
\begin{tikzpicture}[
  node distance=6mm and 10mm,
  startstop/.style={ellipse, draw=black!75, fill=black!8, minimum width=18mm,
    minimum height=7mm, align=center, font=\scriptsize\sffamily},
  action/.style={rectangle, rounded corners=2pt, draw=black!70, fill=black!3,
    minimum height=8mm, text width=42mm, align=center, font=\scriptsize\sffamily},
  decision/.style={diamond, aspect=2.1, draw=black!70, fill=black!4,
    text width=28mm, align=center, font=\scriptsize\sffamily}
]
  \node[startstop] (start) {Inicio};
  \node[action, below=of start] (msg) {Recibir petición del usuario por WebSocket};
  \node[action, below=of msg] (models) {Listar modelos disponibles y recuperar contexto};
  \node[action, below=of models] (architect) {Architect genera especificación de entorno};
  \node[decision, below=of architect] (valid) {¿Es válida y aceptada?};
  \node[action, right=18mm of valid] (revise) {Ajustar petición o regenerar especificación};
  \node[action, below=of valid] (sim) {Cargar modelos y ejecutar simulación};
  \node[action, below=of sim] (tracker) {Tracker registra eventos, trayectorias y episodios};
  \node[action, below=of tracker] (memory) {Persistir observaciones en Knowledge Backbone};
  \node[action, below=of memory] (analyst) {Analyst identifica patrones y genera gráficas};
  \node[action, below=of analyst] (reporter) {Reporter compila informe PDF};
  \node[action, below=of reporter] (ui) {Enviar tarjetas, replay, análisis y enlaces al frontend};
  \node[startstop, below=of ui] (end) {Fin};

  \draw[diagarrow] (start) -- (msg);
  \draw[diagarrow] (msg) -- (models);
  \draw[diagarrow] (models) -- (architect);
  \draw[diagarrow] (architect) -- (valid);
  \draw[diagarrow] (valid) -- node[left, font=\tiny] {sí} (sim);
  \draw[diagarrow] (valid) -- node[above, font=\tiny] {no} (revise);
  \draw[diagarrow] (revise) |- (architect);
  \draw[diagarrow] (sim) -- (tracker);
  \draw[diagarrow] (tracker) -- (memory);
  \draw[diagarrow] (memory) -- (analyst);
  \draw[diagarrow] (analyst) -- (reporter);
  \draw[diagarrow] (reporter) -- (ui);
  \draw[diagarrow] (ui) -- (end);
\end{tikzpicture}}
\caption[Diagrama de actividad UML]{Diagrama de actividad UML de una sesión
completa. El flujo muestra la secuencia canónica del Orchestrator y el punto de
confirmación humana antes de ejecutar la simulación.}
\label{fig:actividad-uml}
\end{figure}

\paragraph{Autocompactación del historial.} El Orchestrator también debe
cuidar su propia ventana de contexto. Una sesión larga acumula conversación,
salidas del Tracker, trazas de decisión y memorias recuperadas; si todo se
mantiene en cada turno, el modelo empieza a trabajar con demasiado ruido.
Para evitarlo, cuando el historial supera un umbral de mensajes o de
caracteres, los turnos antiguos se sustituyen por un resumen de auditoría y
los turnos recientes se conservan intactos. Ese resumen no lo escribe el
LLM: se construye con una plantilla y guarda lo necesario para continuar
(experimento activo, semilla, etapas completadas, peticiones recientes y
herramientas invocadas). Así se libera contexto sin reiniciar la sesión ni
perder la trazabilidad básica de lo ocurrido.

\section{Agentes especializados}

Cuatro agentes especializados completan el pipeline. Todos siguen la
misma forma: reciben un \textit{system prompt} acotado a su rol, una
entrada estructurada del Orchestrator y, cuando procede, un
\texttt{knowledge\_context} pre-recuperado. Devuelven una salida JSON o
Markdown que el Orchestrator interpreta y propaga.

La división en cuatro agentes no responde a una correspondencia directa
con módulos de código, sino a cuatro cambios de representación dentro de
la sesión. El Architect convierte una petición abierta en una
especificación ejecutable; el Tracker convierte una ejecución en
observaciones estructuradas; el Analyst convierte esas observaciones en
interpretación; y el Reporter convierte la interpretación en un documento
revisable. Cada agente opera, por tanto, sobre un tipo de artefacto
distinto y con un criterio de salida propio.

\begin{table}[H]
\begin{center}
\small
\begin{tabular}{|>{\raggedright\arraybackslash}p{2.4cm}|>{\raggedright\arraybackslash}p{4.5cm}|>{\raggedright\arraybackslash}p{5.3cm}|} \hline
\textbf{Agente} & \textbf{Pregunta que resuelve} & \textbf{Razón de separación} \\ \hline
Architect & En qué entorno común se pueden ejecutar los modelos disponibles para el problema descrito. & Requiere traducir lenguaje natural y contexto científico a una especificación formal del dominio de simulación. \\ \hline
Tracker & Qué ocurrió durante la ejecución y qué episodios son relevantes para analizarla. & Requiere separar el registro observacional de la interpretación posterior, conservando datos por paso y por agente. \\ \hline
Analyst & Qué patrones aparecen en las trayectorias y cómo se relacionan con postulados previos. & Requiere combinar métricas, trazas internas, recuperación de conocimiento y comparación entre modelos. \\ \hline
Reporter & Cómo presentar los resultados para que puedan revisarse fuera de la sesión interactiva. & Requiere sintetizar análisis, figuras, métricas y contexto en un documento auto-contenido. \\ \hline
\end{tabular}
\caption[División en agentes]{División en agentes. La tabla resume qué
responsabilidad justifica cada agente especializado.}
\label{tab:razon-agentes-fase2}
\end{center}
\end{table}

Esta frontera es importante porque el Orchestrator no necesita interpretar
texto libre en cada paso: recibe artefactos con una forma esperada y decide
cómo incorporarlos al estado de la sesión.

\paragraph{Architect.} Genera la especificación del entorno de simulación
a partir de la descripción del problema en lenguaje natural y, cuando
existe, del \texttt{knowledge\_\allowbreak context} pre-recuperado por el
Orchestrator. La salida es un JSON que se valida contra el modelo
definido en \texttt{simlab.spec} antes de instanciar el
\texttt{Environment}. El esquema mínimo del JSON contiene tres bloques
de nivel superior: \texttt{grid} con \texttt{width} y \texttt{height};
\texttt{actions}, una lista de acciones con su \texttt{name} y un
\texttt{effect} discriminado (\texttt{MoveEffect}, \texttt{ConsumeEffect}
o \texttt{NoopEffect}); y \texttt{resources}, con los tipos de
recurso, sus propiedades y reglas de aparición. Un ejemplo
representativo es una rejilla \(8 \times 8\) con acciones de movimiento
y consumo, y un único tipo de recurso (comida) que añade energía al
agente que lo consume.

El Architect no decide qué modelo es correcto ni modifica los modelos
cargados. Su responsabilidad es producir un escenario común en el que
puedan ejecutarse en igualdad de condiciones. Para ello recibe también la
lista de modelos disponibles y sus metadatos básicos (nombre, paradigma y
descripción), de forma que la especificación generada no sea una simulación
abstracta cualquiera, sino un entorno compatible con las políticas que se
van a comparar. El Orchestrator conserva la capacidad de presentar esa
propuesta al usuario y de repetir la generación si la especificación no
encaja con la intención de la sesión. El resultado se muestra en la interfaz
como una tarjeta de especificación del entorno; el resto de la interacción se
recoge en la galería del Apéndice~\ref{ap:ui}.

\paragraph{Tracker.} Observa la simulación en bloque, una vez completada,
y produce un JSON estructurado con los eventos por paso, las trayectorias
por agente y los episodios significativos, que la interfaz resume como
tarjetas de trayectoria por agente. Actúa como observador puro:
no escribe en el Knowledge Backbone; esa
responsabilidad recae en el Orchestrator, como se explicó arriba.

La separación entre simulación y observación es deliberada. El motor emite
eventos de bajo nivel (posición, acción, recompensa, percepción y estado
interno) que el Tracker reorganiza en unidades más útiles para el análisis
como trayectoria de cada agente, recursos consumidos, cambios de energía,
eventos críticos, episodios de interés y resumen agregado, evitando que el
Analyst reconstruya la historia de la simulación desde una lista plana de
eventos. Trabaja después de la ejecución, no durante cada paso, porque su
salida depende de patrones que solo se aprecian en la trayectoria completa;
aun así, no introduce interpretación científica: registra y agrupa lo
ocurrido. Esa frontera limpia entre \emph{observación} e \emph{interpretación}
es la que permite que las memorias escritas después distingan hechos de
simulación de conclusiones del análisis.

Los eventos críticos se detectan mediante reglas deterministas antes de que el
LLM interprete la ejecución. El detector marca consumos de recursos, riesgo de
hambre por energía baja, muerte o terminación de un agente, saltos bruscos de
energía, cambios de acción dominante en una ventana deslizante y caídas de
confianza cuando se estrecha la diferencia entre los dos mayores
\textit{Q-values}. De este modo, el Analyst recibe puntos de atención
calculados por código y no episodios inventados por el modelo.

\paragraph{Analyst.} Identifica patrones de comportamiento en los datos del
Tracker, los contrasta con postulados previos y produce un análisis
estructurado de patrones y comparaciones entre modelos
(Figura~\ref{fig:patrones}): es el punto donde la ejecución deja de ser una
colección de eventos y pasa a convertirse en explicación. Su entrada no es la traza bruta,
sino el JSON del Tracker y las ventanas de detalle que el Orchestrator pone a
disposición cuando hacen falta datos concretos; con ellas compara recompensas,
acciones, evolución de variables internas y episodios críticos. Es el único
agente especializado que invoca directamente \texttt{retrieve\_context} como
herramienta; el resto recibe el contexto pre-recuperado. La usa para
matizar la interpretación con conocimiento del backbone, no como sustituto de
los datos observados.

Además del análisis textual, el Analyst genera las visualizaciones que
acompañan al informe: curvas de evolución de variables internas y
recompensas, histogramas de acciones y trazas de decisión por agente
(Figura~\ref{fig:traces}), producidas en \texttt{simlab.charts} a partir de
los datos del Tracker. La
generación no se limita a una única pasada: cada llamada a
\texttt{analyze\_results} puede crear nuevas gráficas, que se acumulan en la
sesión y se muestran en la interfaz. El comportamiento por defecto es generar
un conjunto pequeño que apoye el análisis; si el usuario pide explorar una
variable, intervalo o comparación concreta, el Orchestrator vuelve a invocar
al Analyst con ese foco. Cuando una sesión ejecuta varios modelos en el mismo
entorno (comparativa \emph{multi-modelo}), o se contrasta con ejecuciones
previas de paradigmas distintos (comparativa \emph{multi-paradigma}), produce
gráficas comparadas y un análisis cruzado de las trayectorias
(Figura~\ref{fig:analisis-charts}).

Su salida mantiene una estructura estable para que pueda ser consumida por
el Reporter: resumen de patrones, diferencias entre modelos, métricas
principales, episodios destacados, referencias recuperadas y rutas de las
figuras generadas. La parte generativa del análisis queda así encapsulada en
un artefacto intermedio y el informe final no tiene que
volver a razonar desde los eventos brutos.

\paragraph{Reporter.} Sintetiza los hallazgos del Analyst en un cuerpo
LaTeX que se inserta en la plantilla \texttt{report\_template.tex} y se
compila a PDF mediante \textit{tectonic}. El
contexto del Knowledge Backbone, inyectado por el Orchestrator, le
permite incluir referencias a los paradigmas subyacentes sin tener que
abrir un canal independiente al backbone. El informe incorpora las
gráficas y las tablas de métricas producidas por el Analyst, y se descarga
desde la propia interfaz (Figura~\ref{fig:reporter}).
El nivel de detalle se modula desde la propia configuración del
Reporter según el modelo de lenguaje seleccionado: para modelos
con mayor capacidad de razonamiento se pide una redacción más extensa,
con discusión cualitativa de los patrones, mientras que para modelos
con presupuesto de salida más limitado el informe se reduce a las
secciones esenciales y a las tablas y figuras, manteniendo el documento
auto-contenido.

El Reporter no ejecuta nuevas simulaciones ni consulta directamente el
backbone. Recibe una carpeta de artefactos ya generados y una síntesis
estructurada del Analyst, y su tarea consiste en ordenar esa información
en una pieza legible fuera de la interfaz. La plantilla LaTeX fija la
forma del documento (portada, secciones, tablas, figuras y bibliografía
cuando procede), mientras que el agente se ocupa de redactar la
interpretación y enlazar los resultados con el contexto de la sesión. Así
el informe no es una captura del dashboard, sino un artefacto independiente
que conserva la evidencia principal de la ejecución.

\begin{figure}[H]
\centering
\resizebox{\linewidth}{!}{%
\begin{tikzpicture}[
  actor/.style={rectangle, draw=black!75, rounded corners=2pt, fill=black!5,
    text width=25mm, align=center, font=\scriptsize\sffamily},
  msg/.style={font=\tiny\sffamily, midway, fill=white, inner sep=1pt}
]
  \node[actor] (orch) {Orchestrator};
  \node[actor, right=18mm of orch] (kb) {Knowledge Backbone};
  \node[actor, right=18mm of kb] (analyst) {Analyst + Charts};
  \node[actor, right=18mm of analyst] (reporter) {Reporter};
  \node[actor, right=18mm of reporter] (storage) {MinIO + Frontend};

  \foreach \x in {orch,kb,analyst,reporter,storage}
    \draw[black!35] (\x.south) -- ++(0,-58mm);

  \draw[diagarrow] ([yshift=-9mm]orch.south) -- node[msg] {persistir observaciones} ([yshift=-9mm]kb.south);
  \draw[diagarrow] ([yshift=-18mm]orch.south) -- node[msg] {analyze\_results()} ([yshift=-18mm]analyst.south);
  \draw[diagarrowbi] ([yshift=-27mm]analyst.south) -- node[msg] {retrieve\_context} ([yshift=-27mm]kb.south);
  \draw[diagarrow] ([yshift=-36mm]analyst.south) -- node[msg] {patrones, métricas y figuras} ([yshift=-36mm]orch.south);
  \draw[diagarrow] ([yshift=-45mm]orch.south) -- node[msg] {generate\_report()} ([yshift=-45mm]reporter.south);
  \draw[diagarrow] ([yshift=-54mm]reporter.south) -- node[msg] {PDF + enlaces} ([yshift=-54mm]storage.south);
  \draw[diagarrow] ([yshift=-63mm]storage.south) -- node[msg] {resultado descargable} ([yshift=-63mm]orch.south);
\end{tikzpicture}}
\caption[Diagrama de secuencia de análisis e informe]{Diagrama de secuencia
UML del caso de uso de análisis y generación de informe. El Analyst contrasta
las observaciones con el Knowledge Backbone y el Reporter almacena el PDF como
artefacto descargable.}
\label{fig:secuencia-informe}
\end{figure}

\begin{figure}[H]
\centering
\resizebox{\linewidth}{!}{%
\begin{tikzpicture}[
  actor/.style={rectangle, draw=black!75, rounded corners=2pt, fill=black!5,
    text width=24mm, align=center, font=\scriptsize\sffamily},
  msg/.style={font=\tiny\sffamily, midway, fill=white, inner sep=1pt}
]
  \node[actor] (user) {Usuario + frontend};
  \node[actor, right=17mm of user] (api) {FastAPI};
  \node[actor, right=17mm of api] (orch) {Orchestrator};
  \node[actor, right=17mm of orch] (arch) {Architect};
  \node[actor, right=17mm of arch] (loader) {Model loader};
  \node[actor, right=17mm of loader] (env) {Environment + Tracker};

  \foreach \x in {user,api,orch,arch,loader,env}
    \draw[black!35] (\x.south) -- ++(0,-68mm);

  \draw[diagarrow] ([yshift=-9mm]user.south) -- node[msg] {petición por WebSocket} ([yshift=-9mm]api.south);
  \draw[diagarrow] ([yshift=-17mm]api.south) -- node[msg] {chat()} ([yshift=-17mm]orch.south);
  \draw[diagarrow] ([yshift=-25mm]orch.south) -- node[msg] {generar especificación} ([yshift=-25mm]arch.south);
  \draw[diagarrow] ([yshift=-33mm]arch.south) -- node[msg] {spec validada} ([yshift=-33mm]orch.south);
  \draw[diagarrow] ([yshift=-41mm]orch.south) -- node[msg] {cargar modelos} ([yshift=-41mm]loader.south);
  \draw[diagarrow] ([yshift=-49mm]orch.south) -- node[msg] {ejecutar simulación} ([yshift=-49mm]env.south);
  \draw[diagarrow] ([yshift=-57mm]env.south) -- node[msg] {eventos + JSON observacional} ([yshift=-57mm]orch.south);
  \draw[diagarrow] ([yshift=-65mm]orch.south) -- node[msg] {tarjetas y replay} ([yshift=-65mm]api.south);
  \draw[diagarrow] ([yshift=-73mm]api.south) -- node[msg] {actualización visual} ([yshift=-73mm]user.south);
\end{tikzpicture}}
\caption[Diagrama de secuencia de simulación]{Diagrama de secuencia UML del
caso de uso de ejecución de una simulación. El Orchestrator valida la
especificación, carga modelos de la primera fase y entrega al Tracker la traza
producida por el entorno.}
\label{fig:secuencia-simulacion}
\end{figure}

\section{Carga dinámica de modelos de la primera fase}

El sistema no incorpora estáticamente los modelos generados por la primera
fase. En su lugar, el módulo \texttt{model\_loader} consulta la tabla
\texttt{models} de PostgreSQL (donde la primera fase registra cada
modelo aceptado junto con su paradigma, su formulación, una
descripción y la clave \texttt{s3\_model\_key} del archivo de código en
MinIO) y carga el archivo Python correspondiente bajo demanda en un
directorio temporal por sesión. A continuación importa dinámicamente
las clases que cumplen el contrato esperado. La verificación se hace
por \textit{duck typing}: la segunda fase no impone ninguna clase base ni
dependencia compartida entre ambos sistemas; el \texttt{Protocol}
\texttt{runtime\_checkable} de \texttt{simlab.environment} solo aporta una
comprobación estructural en la carga, y basta con que la clase exponga los
métodos \texttt{decide}, \texttt{update} y \texttt{get\_state}.

Cuando una simulación se ejecuta con semilla, el cargador intenta propagarla
también al modelo. Si el constructor acepta un parámetro \texttt{seed}, se lo
pasa explícitamente; si el módulo generado usa el módulo global
\texttt{random}, lo sustituye durante la carga por un proxy con un generador
sembrado. Esta medida no cambia el contrato público del modelo, pero reduce la
variabilidad accidental y hace que la comparación multi-modelo dependa de la
semilla de la sesión, no de inicializaciones implícitas en el archivo cargado.

Esta decisión tiene tres efectos prácticos. Primero, la incorporación de
un nuevo paradigma no requiere ningún cambio en el código de la segunda
fase. Segundo, el Orchestrator puede presentar al usuario la lista de
modelos disponibles en tiempo real, descubriendo automáticamente los que
la primera fase haya producido entre sesiones. Tercero, una misma
simulación puede instanciar varios modelos sobre el mismo entorno, lo
que habilita comparativas \textit{multi-modelo} sin reconfiguración.

\section{El Knowledge Backbone como capa compartida}

El Knowledge Backbone es la capa de persistencia común a ambas fases del
proyecto. En esta memoria no hace falta repetir toda su definición interna;
para el diseño de la segunda fase interesa, sobre todo, qué consulta el
laboratorio, qué escribe y cómo mantiene la trazabilidad entre piezas.

La capa de almacenamiento combina cuatro componentes con
responsabilidades diferenciadas: PostgreSQL alberga los metadatos, el
historial de experimentos y la tabla
\texttt{simulation\_observations}; MinIO guarda los artefactos binarios;
Qdrant indexa los vectores densos y las representaciones \textit{sparse}
(BM25 nativo) de las memorias; Neo4j mantiene el grafo de paradigmas,
formulaciones y procedencia, que la interfaz expone como un panel navegable
(Figura~\ref{fig:kg}). Los \textit{embeddings} se delegan en Voyage AI y el
\textit{reranking} en ZeroEntropy; ambos se integran en una capa semántica
externa y opcional que requiere las credenciales de los dos proveedores. Si
esa capa no está disponible, no participa: el pipeline principal continúa, pero
sin recuperación semántica, sin \textit{reranking} y sin escritura indexada de
memorias de simulación en Qdrant.

La decisión de repartir la información entre varios almacenes viene de una
necesidad bastante práctica. El sistema conserva metadatos de experimentos y
ejecuciones, modelos registrados por la primera fase, historial conversacional,
artefactos generados y observaciones de simulación reutilizables. No todos esos
datos se consultan igual: unos se filtran por identificador o fecha, otros se
descargan como archivos y otros se recuperan por parecido semántico o por su
relación con paradigmas previos.

\begin{table}[H]
\begin{center}
\small
\setlength{\tabcolsep}{4pt}
\renewcommand{\arraystretch}{1.12}
\begin{tabular}{|>{\raggedright\arraybackslash}p{3.0cm}|>{\raggedright\arraybackslash}p{4.2cm}|>{\raggedright\arraybackslash}p{4.6cm}|} \hline
\textbf{Almacén} & \textbf{Información principal} & \textbf{Motivo de diseño} \\ \hline
PostgreSQL & Experimentos, sesiones, mensajes, modelos registrados y observaciones tipadas. & Consultas exactas, filtrado, trazabilidad relacional y recuperación histórica mediante SQL de solo lectura. \\ \hline
MinIO & Informes PDF, trazas, artefactos binarios y código de modelos aceptados. & Almacenamiento de objetos descargables o de mayor tamaño, separado de las tablas transaccionales. \\ \hline
Qdrant & Memorias densas y representaciones \textit{sparse} BM25 de observaciones y conocimiento recuperable. & Recuperación por similitud semántica o léxica cuando la pregunta no se resuelve bien con claves exactas. \\ \hline
Neo4j & Paradigmas, formulaciones, autores, referencias y relaciones de procedencia. & Navegación por relaciones científicas y contraste entre observaciones y postulados conocidos. \\ \hline
\end{tabular}
\caption[Información almacenada en el Knowledge Backbone]{Reparto de la
información almacenada entre los componentes del Knowledge Backbone.}
\label{tab:almacenamiento-knowledge-backbone}
\end{center}
\end{table}

La coherencia entre almacenes se mantiene mediante identificadores, claves de
artefacto y \textit{UUIDs} compartidos. Por ejemplo, cuando el Reporter genera
un PDF, el Orchestrator guarda la clave del artefacto en el estado de la
sesión, actualiza el experimento y expone enlaces de descarga mediante
\texttt{/api/reports/download} o la herramienta \texttt{get\_report\_links}.

El capítulo anterior explicó por qué se combinan búsqueda semántica,
coincidencia léxica y grafo. Aquí la cuestión práctica es cómo se controla
ese acceso. El Orchestrator pre-recupera contexto para el Architect
y el Reporter, que reciben ese material como \texttt{knowledge\_context}. El
Analyst, en cambio, dispone de \texttt{retrieve\_context} como herramienta
propia para ampliar el contexto cuando una observación concreta lo requiere.
El propio Orchestrator también registra esta herramienta en su catálogo, de
modo que puede responder a preguntas conversacionales sobre paradigmas,
formulaciones u observaciones sin obligar al usuario a navegar manualmente
por los almacenes.

Esa recuperación semántica se complementa con una vía relacional. Cuando el
usuario pregunta por el historial de experimentos, modelos usados, mensajes
previos o memorias persistidas, el Orchestrator puede invocar
\texttt{query\_history}. Esta herramienta traduce la pregunta en lenguaje
natural a una consulta SQL de solo lectura sobre PostgreSQL, valida que solo
use tablas permitidas y devuelve una tabla en \textit{markdown}. Su papel no
es sustituir a Qdrant ni al grafo, sino cubrir preguntas que la similitud de
contenido resuelve peor: recuentos, listados, filtros por estado, cruces con
\texttt{simulation\_observations} o consultas sobre lo que el usuario pidió
en sesiones anteriores.

Existe una segunda ruta NL-SQL, \texttt{query\_experiments}, pensada para
consultas más ricas sobre experimentos. Comparte el mismo principio de SQL de
solo lectura y validación previa, pero puede recuperar artefactos asociados a
los experimentos desde MinIO y pedir una síntesis posterior en lenguaje natural.
En la práctica, \texttt{query\_history} actúa como vía ligera para el chat y
devuelve tablas acotadas, mientras que \texttt{query\_experiments} sirve para
respuestas históricas con más contexto experimental.

La escritura sigue el camino inverso. El Tracker no persiste nada durante su
observación; entrega un JSON al Orchestrator, y este invoca al
\texttt{TrackerMemoryWriter}. El escritor descompone la salida en hechos,
calcula sus \textit{embeddings} en un único \textit{batch} y los inserta en
\texttt{simulation\_observations}, \texttt{memories\_dense} y
\texttt{memories\_sparse}, bajo el \textit{namespace} \texttt{simulation} y
con el mismo \textit{UUID} compartido. Así, una observación de simulación
queda disponible para recuperación posterior sin confundirse con un postulado
procedente de la literatura. En esta fase no se materializa como nodo de Neo4j:
se mantiene como memoria indexada y fila tipada, enlazable por sus metadatos.
El flujo queda, por tanto, concentrado en tres pasos: observación estructurada,
escritura indexada y recuperación posterior desde nuevas sesiones.

\section{Capa de presentación reactiva}

La interfaz no se diseñó como una capa decorativa sobre el pipeline, sino
como parte de la observabilidad del laboratorio. Si el objetivo de la fase
es convertir modelos ejecutables en evidencia revisable, el usuario debe
poder ver cómo se construye esa evidencia: qué agente está trabajando, qué
entorno se ha generado, cómo se mueven los modelos, qué eventos aparecen y
qué análisis se produce al final.

El backend se implementa como una aplicación FastAPI mínima. Tiene un
único canal WebSocket en \texttt{/ws} para la entrada del chat y las
actualizaciones intermedias del pipeline, y un conjunto reducido
de endpoints REST para descargar informes y consultar el Knowledge
Backbone (\texttt{/api/reports/download} y \texttt{/api/knowledge/*}). El
estado de la sesión se mantiene en memoria durante la ejecución y se
persiste en PostgreSQL al cerrarse. Esa persistencia tiene un objetivo
de auditoría y trazabilidad, no de reanudación: el historial de chat se
registra pero no se vuelve a cargar, de modo que el navegador no rehidrata
sesiones pasadas al recargar la página. Ese historial sí puede consultarse
después desde el propio chat mediante \texttt{query\_history}, y también
desde las vistas del Knowledge Backbone cuando el usuario necesita una
inspección manual.

El WebSocket no transporta solo texto final del Orchestrator. También emite
eventos de estado y observabilidad: \texttt{agent\_status} indica cuándo un
agente empieza, termina o vuelve a reposo; \texttt{agent\_tool} muestra qué
herramienta está invocando; las tarjetas intermedias publican la
especificación del entorno, el replay de simulación, el resumen del Tracker
y el análisis con sus gráficas; \texttt{env\_card\_update} refresca datos de
reproducibilidad como semilla y número de pasos; y, cuando el Reporter
genera un PDF, el mensaje final incluye los enlaces descargables. Esta
emisión de eventos forma parte de la observabilidad del sistema: muestra qué
hace cada agente sin esperar a que termine todo el pipeline.

El frontend está construido sobre React 19, Vite y Tailwind CSS 4. La
interfaz principal es un \textit{dashboard} en modo oscuro
(Figura~\ref{fig:dashboard}) que muestra simultáneamente el chat con el
Orchestrator, una visualización animada del entorno con \textit{replay}
paso a paso, tarjetas de datos con métricas y eventos críticos extraídos
por el Tracker, y un panel lateral que refleja el estado de cada agente
mientras ejecuta. La visualización del entorno (Figura~\ref{fig:replay})
reproduce la simulación paso a paso, y al cierre de la sesión el Reporter
genera un informe PDF descargable con las gráficas comparativas de ambos
modelos (Figura~\ref{fig:analisis-charts}). Las capturas de estas tres
vistas se recogen en la galería del Apéndice~\ref{ap:ui}.

El canal WebSocket único busca que cada paso del pipeline aparezca en la
interfaz en cuanto ocurre, sin que el cliente tenga que preguntar por él. La
alternativa de exponer el estado por REST y hacer \textit{polling} desde el
navegador se descartó porque introducía latencia artificial, duplicaba
consultas y convertía una sesión larga en una secuencia de cargas
independientes. El precio es trasladar la sincronización al backend; el
beneficio es que el usuario ve un proceso vivo: el agente piensa, el entorno
se construye y el modelo decide, en lugar de cargas aisladas.

El diseño de la interfaz se documenta mediante capturas reales del sistema
implementado en el Apéndice~\ref{ap:ui}. Se ha optado por capturas en lugar de
\textit{wireframes} porque la interfaz ya estaba materializada al cerrar esta
fase del proyecto; las figuras muestran la organización del \textit{dashboard},
la vista de simulación, el panel del pipeline y las salidas específicas de
Architect, Tracker, Analyst y Reporter.
